% =====================================================================
% BBA INTERNSHIP REPORT TEMPLATE
% Topic: Evaluating MIS Practices in Project Monitoring: A Study on RDRS Bangladesh
% Organization: RDRS Bangladesh, Evaluation Department
% Program: BBA (Major in MIS)
% ---------------------------------------------------------------------
% HOW TO USE:
%  1. Replace all text inside [square brackets] with your own information.
%  2. Compile with pdflatex:  pdflatex main.tex  (run TWICE for TOC/LOT/LOF)
%  3. No exotic packages required - compiles with a standard TeX Live /
%     MiKTeX installation.
% =====================================================================

\documentclass[12pt,a4paper]{report}

% ------------------------- Packages ----------------------------------
\usepackage[margin=1in]{geometry}   % 1-inch margins on A4
\usepackage[T1]{fontenc}            % Better font encoding
\usepackage[utf8]{inputenc}         % UTF-8 input
\usepackage{mathptmx}               % Times-like serif font (professional)
\usepackage{setspace}               % Line spacing control
\usepackage{graphicx}               % Include images: \includegraphics
\usepackage{booktabs}               % Professional tables: \toprule etc.
\usepackage{titlesec}               % Custom chapter/section headings
\usepackage{fancyhdr}               % Custom header/footer
\usepackage{float}                  % [H] placement option for tables/figures
% hyperref MUST be loaded last (clean links, no ugly boxes)
\usepackage[hidelinks]{hyperref}
\hypersetup{
    colorlinks=false,               % No colored boxes; links still clickable
    pdfauthor={[Your Name]},
    pdftitle={Evaluating MIS Practices in Project Monitoring: A Study on RDRS Bangladesh},
    pdfsubject={BBA Internship Report (Major in MIS)},
    pdfkeywords={MIS, Project Monitoring, RDRS Bangladesh, BBA, Internship}
}

% ------------------------- Spacing & Numbering ------------------------
\onehalfspacing                     % 1.5 line spacing for body text
\setcounter{secnumdepth}{2}         % Number \chapter, \section, \subsection
\setcounter{tocdepth}{2}            % Show subsections in the TOC

% ------------------------- Heading Styles -----------------------------
% Chapter: bold, large, with horizontal rule underneath
\titleformat{\chapter}[display]
    {\normalfont\bfseries\Large}
    {\chaptername\ \thechapter}
    {10pt}
    {\LARGE}
    [\vspace{8pt}\rule{\textwidth}{1pt}]

% Section / Subsection: bold, clean sizes
\titleformat{\section}
    {\normalfont\large\bfseries}{\thesection}{1em}{}
\titleformat{\subsection}
    {\normalfont\normalsize\bfseries}{\thesubsection}{1em}{}

% Extra spacing around headings
\titlespacing*{\chapter}{0pt}{-20pt}{20pt}
\titlespacing*{\section}{0pt}{14pt}{8pt}
\titlespacing*{\subsection}{0pt}{10pt}{6pt}

% ------------------------- Header / Footer ----------------------------
\setlength{\headheight}{14.5pt}     % Avoid fancyhdr warnings
\pagestyle{fancy}
\fancyhf{}                          % Clear defaults
% TODO: Replace header text if needed, or leave empty for a minimal look
\fancyhead[L]{\footnotesize BBA Internship Report}          % Left header
\fancyhead[R]{\footnotesize RDRS Bangladesh}                % Right header
\fancyfoot[C]{\thepage}             % Page number centered in footer
\renewcommand{\headrulewidth}{0.4pt}
\renewcommand{\footrulewidth}{0pt}
% Keep chapter opening pages with centered page numbers too
\fancypagestyle{plain}{%
    \fancyhf{}%
    \fancyfoot[C]{\thepage}%
    \renewcommand{\headrulewidth}{0pt}%
}

% =====================================================================
\begin{document}

% ======================= COVER / TITLE PAGE ==========================
% TODO: Replace ALL placeholder text below with your actual details.
\begin{titlepage}
    \centering

    % --- University Name ---
    {\Large\bfseries [Your University Name] \par}
    \vspace{4pt}
    {\large Faculty of Business Administration \par}
    {\normalsize Department of Management Information Systems \par}

    \vspace{12pt}
    \rule{\textwidth}{1.5pt}
    \vspace{12pt}

    % --- Report Title ---
    {\LARGE\bfseries Evaluating MIS Practices in\\[6pt] Project Monitoring:\\[6pt]
        A Study on RDRS Bangladesh \par}
    \vspace{12pt}
    {\large\itshape An Internship Report on the Evaluation Department,
        RDRS Bangladesh \par}

    \vspace{16pt}

    % --- Submitted By / Submitted To boxes ---
    \begin{minipage}[t]{0.46\textwidth}
        \raggedright
        \textbf{Submitted By:}\\[4pt]
        [Your Full Name]\\           % TODO: your name
        ID: [XX-XXXXX-X]\\           % TODO: your student ID
        Program: BBA (Major in MIS)\\% Keep as is, or update batch/section
        Batch: [XXth Batch]\\        % TODO: your batch
        Semester: [Spring/Fall 20XX] % TODO: semester
    \end{minipage}
    \hfill
    \begin{minipage}[t]{0.46\textwidth}
        \raggedright
        \textbf{Submitted To:}\\[4pt]
        [Supervisor's Full Name]\\   % TODO: academic supervisor
        [Designation, e.g. Assistant Professor]\\ % TODO: designation
        Department of MIS\\          % TODO: adjust if needed
        [Your University Name]\\
        \vspace{6pt}
        \textbf{Organization Supervisor:}\\ % TODO: RDRS supervisor
        [Name, Designation]\\
        Evaluation Department, RDRS Bangladesh
    \end{minipage}

    \vfill

    % --- Date ---
    {\large Date of Submission: [DD Month YYYY] \par} % TODO: submission date
    \vspace{8pt}
    \rule{\textwidth}{1.5pt}
    \vspace{6pt}
    {\small In partial fulfillment of the requirements for the degree of
        Bachelor of Business Administration (BBA), Major in Management
        Information Systems \par}
\end{titlepage}
% ===================== END COVER PAGE =================================

% Front matter uses roman numerals (i, ii, iii, ...)
\pagenumbering{roman}

% ===================== LETTER OF TRANSMITTAL ==========================
% Unnumbered chapter (does not get a chapter number, but appears in TOC).
\chapter*{Letter of Transmittal}
\addcontentsline{toc}{chapter}{Letter of Transmittal}

% TODO: Replace date, names, and addresses.
\noindent
[DD Month YYYY] \\[6pt]
To\\
[Supervisor's Full Name]\\
[Designation]\\
Department of Management Information Systems\\
[Your University Name]\\
[University Address]\\[12pt]

\noindent
\textbf{Subject: Submission of Internship Report on ``Evaluating MIS
Practices in Project Monitoring: A Study on RDRS Bangladesh.''}\\[12pt]

\noindent
Dear Sir/Madam,\\[6pt]

With due respect, I, [Your Full Name], ID: [XX-XXXXX-X], a student of the
BBA program (Major in MIS), [Your University Name], have successfully
completed my internship at the Evaluation Department of RDRS Bangladesh,
Rangpur, from [Start Date] to [End Date]. As a requirement of the BBA
program, I am pleased to submit my internship report entitled
``Evaluating MIS Practices in Project Monitoring: A Study on RDRS
Bangladesh.''

During my internship, I observed how project data is collected from field
offices, validated at the Evaluation Department, and used for monitoring
reports and management decisions. This report analyzes the existing MIS
tools, data flow, reporting practices, and challenges, and proposes
practical recommendations.

I have prepared this report with sincerity and to the best of my
knowledge. I would be grateful for your kind evaluation and feedback.

\vspace{12pt}
\noindent
Sincerely yours,\\[24pt]
\rule{2.5in}{0.5pt}\\
\textbf{[Your Full Name]}\\
ID: [XX-XXXXX-X]\\
BBA (Major in MIS), [Your University Name]\\
Contact: [Phone / Email] % TODO: your contact
\newpage

% ========================= ACKNOWLEDGEMENT =============================
\chapter*{Acknowledgement}
\addcontentsline{toc}{chapter}{Acknowledgement}

% TODO: Personalize with real names.
First and foremost, I express my gratitude to Almighty Allah for giving me
the strength to complete this internship and report successfully.

I am sincerely thankful to my academic supervisor, [Supervisor's Name,
Designation], Department of MIS, [Your University Name], for his/her
valuable guidance, constructive feedback, and continuous encouragement
throughout this work.

I am equally grateful to my organization supervisor, [Name, Designation],
Evaluation Department, RDRS Bangladesh, and to all officers and field
staff who spared their busy time to explain the monitoring system,
share data collection formats, and answer my interview questions. Without
their cooperation, this study on MIS practices in project monitoring
would not have been possible.

Finally, I thank my family, friends, and fellow interns for their moral
support during the internship period.

\vspace{12pt}
\noindent
[Your Full Name]\\
BBA (Major in MIS)\\
[Your University Name]
\newpage

% ======================= EXECUTIVE SUMMARY =============================
\chapter*{Executive Summary}
\addcontentsline{toc}{chapter}{Executive Summary}

% TODO: Rewrite in your own words after completing Chapters 4-6 (~300-400 words).
This internship report evaluates the Management Information System (MIS)
practices used for project monitoring at RDRS Bangladesh, with special
focus on the Evaluation Department. The study was conducted during a
[three-month] internship from [Start Date] to [End Date] at the RDRS
head office in Rangpur.

The broad objective of the study was to assess how effectively MIS tools
and processes support timely, accurate, and decision-oriented project
monitoring. Data were collected from both primary sources (direct
observation, informal discussions, and semi-structured interviews with
[X] Evaluation Department staff and field officers) and secondary
sources (project documents, monitoring formats, progress reports, and
the RDRS website and publications).

The findings show that RDRS Bangladesh maintains a structured,
largely Excel- and paper-assisted MIS supported by periodic field
verification. The system performs well in routine data aggregation and
donor reporting, but faces challenges including duplicate data entry,
delayed field-level reporting, limited dashboard-based visualization,
and uneven digital literacy among field staff. The study identifies
specific gaps in data validation, real-time tracking, and feedback
loops between the Evaluation Department and project teams.

Based on these findings, the report recommends a phased digitization of
field data collection, standardized indicator definitions, automated
validation checks, a centralized monitoring dashboard, and regular MIS
training for field and head-office staff. Strengthening these areas
would improve data quality, reduce reporting delays, and enhance the
use of monitoring information in management decisions.

In conclusion, RDRS Bangladesh has a committed evaluation culture and a
functional MIS foundation. With targeted improvements in automation,
integration, and capacity building, its MIS can become a more
proactive, real-time project monitoring system.
\newpage

% ============ TABLE OF CONTENTS / LISTS OF TABLES & FIGURES ============
\tableofcontents
\newpage
\listoftables
\listoffigures
\newpage

% Main matter uses arabic numerals (1, 2, 3, ...)
\pagenumbering{arabic}

% ========================= CHAPTER 1 ===================================
\chapter{Introduction}

\section{Background of the Study}
% TODO: Rewrite with your own context and dates.
In the contemporary development sector, timely and reliable information is
critical for effective project management. Management Information Systems
(MIS) provide the tools and processes through which organizations
collect, process, store, and disseminate project data for monitoring and
decision-making. For non-governmental organizations (NGOs) like RDRS
Bangladesh, which implement multiple donor-funded projects across
northern Bangladesh, a well-functioning MIS is essential to track
progress against indicators, ensure accountability, and demonstrate
results to stakeholders.

RDRS Bangladesh, formerly the Rangpur Dinajpur Rural Service, has evolved
from a relief-oriented program into a leading national development
organization. Its Evaluation Department plays a central role in
monitoring project performance, verifying field-level data, and
producing progress reports. This report examines how MIS practices in
this department support project monitoring, what works well, and where
improvements are needed. The study connects classroom concepts from MIS,
database management, and project monitoring with real organizational
practices observed during the internship.

\section{Objectives of the Study}
The study is guided by one broad objective and several specific
objectives.

\subsection{Broad Objective}
% TODO: Keep or refine to match your approved proposal.
The broad objective of this study is to evaluate the existing MIS
practices in project monitoring at RDRS Bangladesh and to suggest
improvements for more effective, timely, and data-driven monitoring.

\subsection{Specific Objectives}
% TODO: Adjust the bullet list to match your approved objectives.
\begin{itemize}
    \item To understand the current MIS tools, data flow, and reporting
        process used by the Evaluation Department.
    \item To assess the effectiveness of MIS practices in ensuring data
        accuracy, timeliness, and usability for project monitoring.
    \item To identify the key challenges and gaps in data collection,
        validation, storage, and reporting.
    \item To recommend practical measures for strengthening MIS-based
        project monitoring at RDRS Bangladesh.
\end{itemize}

\section{Scope of the Study}
% TODO: Adjust location, duration, and coverage.
This study focuses on the Evaluation Department at the RDRS Bangladesh
head office in Rangpur and its linkage with selected field offices. It
covers the MIS practices used for routine project monitoring --- data
collection formats, data entry and validation, storage, analysis, and
reporting --- during the internship period from [Start Date] to
[End Date]. The study does not attempt a full IT audit or a financial
evaluation; it concentrates on the information flow and monitoring use
of MIS from a management perspective.

\section{Methodology}
This section describes the sources of data and the methods used to
collect and analyze information.

\subsection{Sources of Data}
Both primary and secondary data were used:
\begin{itemize}
    \item \textbf{Primary sources:} Direct observation of daily MIS
        activities, informal discussions with Evaluation Department
        staff, and semi-structured interviews with [X] respondents
        including monitoring officers and field facilitators.
    \item \textbf{Secondary sources:} Project proposals, logical
        frameworks, monitoring formats, monthly/quarterly progress
        reports, RDRS annual reports, and the official RDRS website
        and published documents.
\end{itemize}

\subsection{Data Collection Method}
% TODO: Update sample size and tools.
Data were collected through (i) participant observation during regular
office work, (ii) a semi-structured interview checklist (see Appendix),
and (iii) review of existing monitoring documents and report templates.
Qualitative responses were summarized thematically around data quality,
timeliness, usability, and system challenges. Simple frequency counts
were used where applicable to highlight patterns in the findings
presented in Chapter~4.

\section{Limitations of the Study}
% TODO: Be honest about your actual constraints.
This study has several limitations. First, the internship duration of
[three months] limited the scope of field visits; most observations are
based on the head office and [one/two] field offices. Second, some
project data are confidential and could not be disclosed in this report;
aggregated and anonymized examples are used instead. Third, the study
relies partly on respondents' perceptions, which may involve recall or
response bias. Despite these limitations, the findings provide a useful
overview of MIS practices in project monitoring at RDRS Bangladesh.

% ========================= CHAPTER 2 ===================================
\chapter{Organizational Overview}

\section{About RDRS Bangladesh}
% TODO: Verify facts (founding year, coverage, sectors) from RDRS sources.
RDRS Bangladesh is a leading national development organization working to
reduce poverty and empower rural communities, particularly in the
northwestern region of Bangladesh. Established in [1972 as a field
program and later registered as a national NGO], RDRS works in areas
such as education, health, agriculture, climate resilience, governance,
and microfinance. The organization partners with government agencies,
international donors, and local communities to implement integrated
development projects. During the internship, the author observed a strong
field presence coordinated through regional and local offices reporting
to the Rangpur head office.

\section{Mission, Vision, and Objectives}
% TODO: Copy the official statements from RDRS publications/website.
\begin{itemize}
    \item \textbf{Vision:} [e.g. A just and equitable society where the
        poor exercise their rights and live with dignity.] % TODO: verify
    \item \textbf{Mission:} [e.g. To empower rural communities through
        sustainable development programs, capacity building, and policy
        advocacy.] % TODO: verify
    \item \textbf{Objectives:} (i) Strengthen livelihoods and food
        security; (ii) improve access to quality education and health
        services; (iii) build resilience to climate change; (iv) promote
        good governance and social inclusion; and (v) ensure accountable,
        evidence-based project management through effective monitoring
        and evaluation.
\end{itemize}

\section{Organizational Structure}
RDRS Bangladesh is headed by the Executive Director, supported by
directors of program, finance, and administration divisions. The
Evaluation Department functions as a cross-cutting support unit,
working with sectoral programs and field management. A simplified
organogram is shown in Figure~\ref{fig:orgchart}.

% TODO: Replace the placeholder box with your org chart image:
%   \includegraphics[width=0.9\textwidth]{images/org-chart.png}
\begin{figure}[H]
    \centering
    \fbox{\parbox[c][2.8in][c]{4.2in}{%
        \centering \textit{[Insert Organizational Chart of RDRS Bangladesh here]}\\
        \small Executive Director $\rightarrow$ Directors $\rightarrow$
        Evaluation Department + Programs + Field Offices}}
    \caption{Simplified Organizational Structure of RDRS Bangladesh (placeholder)}
    \label{fig:orgchart}
\end{figure}

\section{The Evaluation Department}
% TODO: Confirm staffing and responsibilities with your supervisor.
The Evaluation Department is responsible for designing monitoring
frameworks, developing data collection tools, verifying field data,
analyzing progress against indicators, and preparing monitoring reports
for management and donors. Key functions observed during the internship
include: (i) reviewing monthly field reports, (ii) conducting spot
checks and joint field verification, (iii) maintaining project
databases and filing systems, and (iv) providing feedback to project
teams for corrective action. The department typically uses a combination
of paper formats, spreadsheets, and shared files, with gradual adoption
of digital data collection tools.

\section{Ongoing Projects Overview}
% TODO: List 3-5 projects you actually observed (use public information only).
RDRS Bangladesh implements multiple projects at any given time. Examples
observed or reviewed during the internship include:
\begin{itemize}
    \item \textbf{[Project A --- e.g. Climate-Resilient Livelihoods]:}
        Focus on [adaptive agriculture and income diversification];
        donor: [Donor name]. % TODO: replace
    \item \textbf{[Project B --- e.g. Quality Primary Education]:}
        Focus on [school support and community engagement];
        donor: [Donor name]. % TODO: replace
    \item \textbf{[Project C --- e.g. Community Health and Nutrition]:}
        Focus on [maternal health and nutrition awareness];
        donor: [Donor name]. % TODO: replace
\end{itemize}
Each project has its own results framework, but all report through a
common monitoring cycle coordinated by the Evaluation Department.

% ========================= CHAPTER 3 ===================================
\chapter{Literature Review}
% TODO: Expand with 4-6 textbook/journal sources from your course list.

\section{Concept of MIS}
A Management Information System (MIS) is an organized system for
collecting, processing, storing, and distributing information to support
decision-making and control in an organization \cite{laudon2022}. In the
context of development projects, MIS transforms raw field data --- such
as beneficiary counts, training records, and input distribution --- into
structured monitoring information such as progress dashboards, variance
reports, and trend analyses.

\section{MIS in Project Monitoring}
Project monitoring is the continuous process of tracking implementation
against plans and indicators. Literature emphasizes that effective
monitoring depends on clearly defined indicators, standardized data
collection, validation mechanisms, and timely feedback loops
\cite{naidoo2011}. MIS supports each of these stages by enforcing
consistent formats, automating aggregation, and making information
accessible to managers for corrective decisions.

\section{MIS in NGOs in Bangladesh}
Studies on Bangladeshi NGOs note a transition from paper-based reporting
toward spreadsheet- and mobile-based systems, with persistent challenges
in data quality, connectivity, and staff capacity \cite{nahmias2020}.
Organizations with dedicated M\&E units and clear data-flow protocols
tend to produce more reliable monitoring reports, while fragmented tools
and duplicate entry increase errors and delays.

\section{Theoretical Framework and Research Gap}
This report adopts a simple input--process--output view of MIS for
monitoring: field data (input) are validated and processed by the
Evaluation Department (process) to produce monitoring reports and
decisions (output). While existing literature discusses NGO monitoring
systems generally, few studies examine the day-to-day MIS practices of
the Evaluation Department of RDRS Bangladesh specifically. This report
addresses that gap through direct observation and staff interviews.

% ========================= CHAPTER 4 ===================================
\chapter{Analysis and Findings}

\section{Overview of the Current MIS Workflow}
% TODO: Describe the actual data flow you observed.
Field staff record activity data in prescribed paper or digital formats
and submit them monthly to theEvaluation Department through field
supervisors. The department checks completeness, enters or consolidates
data into spreadsheets, resolves queries with field offices, and
prepares monthly and quarterly summaries for project managers and
donors. Spot verification visits are conducted to validate reported
figures. This workflow is functional but involves repeated manual
handling at several stages.

\section{Effectiveness of MIS Practices}
Table~\ref{tab:findings} summarizes the author's assessment of key MIS
dimensions based on observation and interview responses
(\(n =\) [X] respondents). % TODO: update sample size and ratings.
Scores use a simple 5-point scale (1 = Very Weak, 5 = Very Strong) for
illustration.

% TODO: Replace the example numbers with your actual survey/assessment results.
\begin{table}[H]
    \centering
    \caption{Assessment of MIS Practices in Project Monitoring (example format)}
    \label{tab:findings}
    \begin{tabular}{lccc}
        \toprule
        \textbf{MIS Dimension} & \textbf{Mean Score} & \textbf{Rank} & \textbf{Remarks} \\
        \midrule
        Data accuracy \& validation   & 3.8 & 1 & Spot checks help; entry errors remain \\
        Timeliness of reporting       & 3.2 & 3 & Delays from field submission \\
        Ease of data entry            & 3.5 & 2 & Formats clear; some duplication \\
        Dashboard \& visualization    & 2.6 & 5 & Mostly tables; few charts \\
        Use in decision-making        & 2.9 & 4 & Reports used; feedback loop slow \\
        \bottomrule
    \end{tabular}
\end{table}

The pattern suggests that data accuracy is relatively stronger due to
manual verification, while visualization and feedback use are weaker.
Respondents appreciated the clarity of monitoring formats but reported
that duplicate entry across multiple spreadsheets consumes significant
time.

\section{Visualization and Reporting Practice}
Figure~\ref{fig:reporting} illustrates a typical reporting pattern
observed during the internship: most monitoring outputs are tabular
progress tables, with limited use of charts or dashboards for trend
analysis. % TODO: Replace with your own chart (e.g. export from Excel).

% TODO: Replace the placeholder box with your own figure:
%   \includegraphics[width=0.85\textwidth]{images/reporting-example.png}
\begin{figure}[H]
    \centering
    \fbox{\parbox[c][2.5in][c]{4.2in}{%
        \centering \textit{[Insert example monitoring chart or dashboard screenshot here]}\\
        \small e.g. Monthly beneficiary coverage vs.\ target (bar chart)}}
    \caption{Example of Current Monitoring Report Visualization (placeholder)}
    \label{fig:reporting}
\end{figure}

\section{Key Challenges Identified}
% TODO: Keep only the challenges you actually found evidence for.
\begin{itemize}
    \item \textbf{Duplicate data entry:} The same field data are entered
        in multiple files for different projects and donors.
    \item \textbf{Delayed field reporting:} Late submission from remote
        offices compresses the time available for validation.
    \item \textbf{Limited automation:} Few validation rules or automated
        checks; errors are caught manually.
    \item \textbf{Weak visualization:} Managers receive long tables with
        limited summary charts or exception highlights.
    \item \textbf{Capacity gaps:} Uneven Excel and digital-tool skills
        among field staff slow down adoption of new formats.
\end{itemize}

% ========================= CHAPTER 5 ===================================
\chapter{Recommendations}
% TODO: Ensure each recommendation links to a finding in Chapter 4.
Based on the analysis above, the following recommendations are proposed
for the Evaluation Department and management of RDRS Bangladesh:

\begin{enumerate}
    \item \textbf{Digitize field data collection in phases.} Introduce
        simple mobile-based forms (with offline capability) for routine
        monthly reporting to reduce paper handling and transcription
        errors. Start with one pilot project before scaling up.
    \item \textbf{Standardize indicators and formats.} Publish a single
        indicator dictionary with clear definitions, disaggregation
        rules, and one approved format per data type to eliminate
        parallel versions.
    \item \textbf{Add automated validation checks.} Build range checks,
        mandatory-field rules, and duplicate warnings into spreadsheet
        templates or data entry forms to catch errors at the source.
    \item \textbf{Create a centralized monitoring dashboard.} Develop a
        simple dashboard (e.g. monthly coverage vs.\ target, overdue
        reports, data-quality flags) for project managers using existing
        tools before investing in larger software.
    \item \textbf{Strengthen the feedback loop.} Share validated
        monitoring summaries with field offices within [X] working days,
        with clear action points, so reporting is seen as useful rather
        than purely compliance-driven.
    \item \textbf{Invest in regular MIS training.} Conduct short,
        hands-on training on data entry, basic Excel functions, and data
        ethics for field and head-office staff at least twice a year.
\end{enumerate}

% ========================= CHAPTER 6 ===================================
\chapter{Conclusion}
% TODO: Summarize without introducing major new points.
This report evaluated the MIS practices in project monitoring at the
Evaluation Department of RDRS Bangladesh. The study found a committed
team, clear reporting formats, and a functioning verification system
that together support credible routine monitoring and donor reporting.
At the same time, manual processes, duplicate entry, reporting delays,
and limited visualization constrain the timeliness and decision-use of
monitoring information.

The recommendations --- phased digitization, standardization, automated
validation, a centralized dashboard, faster feedback, and regular
training --- are practical steps that build on existing strengths
rather than requiring an immediate large-scale system overhaul. With
these improvements, the Evaluation Department can move from a largely
retrospective reporting role toward a more real-time, analytical
monitoring function.

As an intern, the author gained hands-on understanding of how MIS
concepts such as data quality, system workflow, and information use
apply in a development organization. The experience bridged academic
learning in the BBA (MIS) program with the realities of field-level
project monitoring, and deepened appreciation for evidence-based
management in the NGO sector.

% ========================= REFERENCES ==================================
% Manually formatted in APA style (7th ed.) as an example.
% TODO: Replace with the sources you actually cited. Keep alphabetical order.
\begin{thebibliography}{99}

\bibitem{laudon2022}
Laudon, K. C., \& Laudon, J. P. (2022). \textit{Management information
systems: Managing the digital firm} (17th ed.). Pearson.

\bibitem{naidoo2011}
Naidoo, I. (2011). \textit{Monitoring and evaluation practices and
challenges of non-governmental organisations}. [Publisher/Institution].
% TODO: complete with actual source details.

\bibitem{nahmias2020}
Nahmias, P., \& Baumann, H. (2020). Digital monitoring systems in
development programmes: Lessons from South Asia. \textit{Journal of
Development Informatics}, \textit{4}(2), 45--62.
% TODO: replace with a real article you consulted, or delete this entry.

\bibitem{rdrs2024}
RDRS Bangladesh. (2024). \textit{Annual report 2023--2024}. RDRS
Bangladesh, Rangpur.
% TODO: verify year/title and add URL if accessed online.

\end{thebibliography}

% ========================= APPENDIX ====================================
\appendix
\chapter{Appendix A: Interview Checklist and Sample Formats}
% TODO: Insert your actual approved instruments here.

\section*{A.1 Semi-Structured Interview Checklist (Sample)}
\begin{enumerate}
    \item What is your role in the project monitoring and reporting process?
    \item Which data collection formats/tools do you use most frequently?
    \item How is field-level data verified before final reporting?
    \item What challenges do you face in data entry, validation, or timely submission?
    \item How are monitoring reports used by management for decision-making?
    \item What improvements to the current MIS would you suggest?
\end{enumerate}

\section*{A.2 Sample Monitoring Format (Placeholder)}
% TODO: Paste a blank/anonymized version of one real format (remove sensitive data).
\begin{center}
\fbox{\parbox[c][2.2in][c]{4.5in}{%
    \centering \textit{[Attach a blank monthly field-report format here]}\\
    \small Remember to anonymize beneficiary and staff identifiers.}}
\end{center}

\end{document}
